% Einheit 5 von 5 – Prozentuale Veränderung (bank/prozentrechnung/e5.jsonl, zusammenbau v0.1)
%% TODO Zweigzeile Teil 1: was hier gelernt wird (Satz in Schülersprache aus dem Einheitstitel) – Einheit 5
\einheitenkopf[e5]{Einheit 5 von 5 · Prozentuale Veränderung}
\zweigzeile{neu in diesem Jahr, je nach Buch bis Klasse 8 · P10 oft}
\setcounter{aufgabe}{22}

%% TODO Titel: Ich-kann-Satz fehlt in der Bank (Platzhalter: Kettenname) – Nr. 23
%% TODO Anweisung: ein Satz über der ersten Teilaufgabe fehlt in der Bank – Nr. 23
\begin{aufgabe}{Welcher Wert ist der alte?}
\begin{teile}
\teil Der Preis stieg von $2{,}40\,€$ auf $2{,}70\,€$. – Welcher Wert ist der alte?
\end{teile}
\end{aufgabe}

%% TODO Titel: Ich-kann-Satz fehlt in der Bank (Platzhalter: Kettenname) – Nr. 24
%% TODO Anweisung: ein Satz über der ersten Teilaufgabe fehlt in der Bank – Nr. 24
\begin{aufgabe}{Veränderung}
\begin{teile}
\teil um $20\,\%$ gesenkt – was bedeutet dasselbe? Kreuze an.\\ \kreuz{auf $20\,\%$ gesenkt}\\ \kreuz{auf $80\,\%$ gesenkt}\\ \kreuz{um $80\,\%$ gesenkt}
\end{teile}
%% TODO Darstellung für schwach fehlt in der Bank (prozentrechnung-e5-k2-s1-v1; Abschnitt „Für schwache Schüler“)
\swz{$80\,€$ um $10\,\%$ erhöht – neuer Preis?}{}
%% TODO Darstellung für schwach fehlt in der Bank (prozentrechnung-e5-k2-s1-v2; Abschnitt „Für schwache Schüler“)
\swz{$60$ kg um $25\,\%$ gesenkt – neues Gewicht?}{}
%% TODO Darstellung für schwach fehlt in der Bank (prozentrechnung-e5-k2-s1-v3; Abschnitt „Für schwache Schüler“)
\swz{$400$ Einwohner, $5\,\%$ mehr – wie viele jetzt?}{}
%% TODO Darstellung für schwach fehlt in der Bank (prozentrechnung-e5-k2-s1-v4; Abschnitt „Für schwache Schüler“)
\swz{$50\,€$ um $20\,\%$ gesenkt – neuer Preis?}{}
%% TODO Darstellung für schwach fehlt in der Bank (prozentrechnung-e5-k2-s1-v5; Abschnitt „Für schwache Schüler“)
\swz{$90$ cm um $50\,\%$ verlängert – neue Länge?}{}
\swfrage{Was bleibt gleich, was ändert sich?}
%% TODO Darstellung für schwach fehlt in der Bank (prozentrechnung-e5-k2-s2-v1; Abschnitt „Für schwache Schüler“)
\swz{$45\,€$ um $20\,\%$ erhöht – mit Faktor}{}
%% TODO Darstellung für schwach fehlt in der Bank (prozentrechnung-e5-k2-s3-v1; Abschnitt „Für schwache Schüler“)
\swz{von $40$ auf $50$ – um wie viel Prozent mehr?}{}
%% TODO Darstellung für schwach fehlt in der Bank (prozentrechnung-e5-k2-s4-v1; Abschnitt „Für schwache Schüler“)
\swz{Nach $20\,\%$ Rabatt kostet ein Rad $360\,€$ – alter Preis?}{}
%% TODO Darstellung für schwach fehlt in der Bank (prozentrechnung-e5-k2-s5-v1; Abschnitt „Für schwache Schüler“)
\swz{netto (ohne Steuer) $200\,€$, $19\,\%$ Mehrwertsteuer – brutto?}{}
%% TODO Darstellung für schwach fehlt in der Bank (prozentrechnung-e5-k2-s6-v1; Abschnitt „Für schwache Schüler“)
\swz[3]{Wahlbeteiligung $58\,\%$, später $64\,\%$ – wie viele Prozentpunkte mehr? Wie viel Prozent mehr? Runde auf eine Stelle.}{}
%% TODO Darstellung für schwach fehlt in der Bank (prozentrechnung-e5-k2-s7-v1; Abschnitt „Für schwache Schüler“)
\swz{Eine Straße steigt auf $200$ m waagerechter Strecke um $14$ m – Steigung in Prozent?}{}
\swz[3]{Ein Liter Milch kostete $1{,}15\,€$, jetzt kostet er $1{,}29\,€$. Um wie viel Prozent ist der Preis gestiegen? Runde auf eine Stelle. \hfill (P10 2026 FOR)}{}
\swz[3]{Ein Busticket kostet $2{,}40\,€$ und wird um $25\,\%$ teurer. Wie viel kostet es dann? \hfill (P10 2024 OS)}{}
\swz[3]{Kletterhalle: Erwachsene $14\,€$, Jugendliche (bis $17$ Jahre) $9\,€$. Familie Demir: Mutter, Vater, Tochter ($15$ Jahre), Sohn ($12$ Jahre). Montags gibt es $15\,\%$ Rabatt. Wie viel zahlt die Familie am Montag? \hfill (P10 2015 OS)}{}
\swz[3]{Für Radwege gilt: Steigung höchstens $5\,\%$. Ergänze: „$5\,\%$ Steigung heißt: auf \leerfeld[m] waagerechter Strecke \leerfeld[m] Höhenunterschied.“ Herr Lenz sagt: „Ein Radweg, der auf $240$ m waagerechter Strecke um $15$ m ansteigt, ist zu steil.“ Hat er recht? Rechne. \hfill (P10 2025 OS)}{}
\end{aufgabe}

%% TODO Titel: Ich-kann-Satz fehlt in der Bank (Platzhalter: Kettenname) – Nr. 25
%% TODO Anweisung: ein Satz über der ersten Teilaufgabe fehlt in der Bank – Nr. 25
\begin{aufgabe}{Veränderung – Fehler finden, Begründen, Anwendung}
\swz[3]{Preis von $60\,€$ auf $75\,€$ – Steigerung in Prozent? Ida: \rechnung{15 : 75 &= 0{,}2 = 20\,\%} Fehler? Wie heißt es richtig?}{}
\swfrage{Warum bezieht man eine Preiserhöhung auf den alten Preis?}
\swz[3]{Zwei Angebote für ein Fahrrad zu $600\,€$: A gibt $10\,\%$ Rabatt und auf den neuen Preis nochmals $10\,\%$; B gibt einmal $20\,\%$ Rabatt. Welches Angebot ist günstiger?}{}
\end{aufgabe}

\uebersichtskasten{Erhöhung um p \%: neuer Wert = alter Wert · (1 + p/100); Senkung: · (1 − p/100). \\ 250 € um 12 \% erhöht: 250 · 1,12 = 280 € \quad 250 € um 12 \% gesenkt: 250 · 0,88 = 220 € \\ Veränderung in Prozent: (neu − alt) : alt. \quad von 60 auf 69: 9 : 60 = 0,15 = 15 \% \\ Prozentpunkte: Unterschied zweier Prozentangaben. \quad von 30 \% auf 33 \%: 3 Prozentpunkte (das sind 10 \% mehr) \\ Steigung in Prozent: Höhe : waagerechte Strecke. \quad 8 m auf 100 m = 8 \%}
